% Zdroj: PDF priklady-leto-2024.pdf, fyzická str. 24 (zadání), obrázek str. 25. Značka: specializace UI-SU, UI-ZPJ. Zařazeno: S3.
\section{Metoda podpůrných vektorů (SVM)}
\szzotazka{léto 2024}{30}{Metoda podpůrných vektorů (SVM)}

\noindent\textbf{Zadání.}\par
Uvažujme data se dvěmi příznaky, která chceme klasifikovat do dvou tříd. Data pro třídu A jsou
na pozicích $(1,1)$, $(2,2)$ a $(2,0)$, data pro třídu B jsou na pozicích $(0,0)$, $(1,0)$ a $(0,1)$.
Data jsou také zobrazena na obrázku níže. Třída A je označena jako křížek ($\times$), třída B jako
kolečko ($\bullet$).

\begin{center}\includegraphics[width=0.6\linewidth]{svm.png}\end{center}

\begin{enumerate}
  \item Popište, jak fungují a jak se trénují lineární support vector machines (metoda podpůrných
    vektorů) v případě, že předpokládáme, že data jsou lineárně separabilní.
  \item Na základě dat výše zakreslete přibližnou rozhodovací hranici linárního SVM, která odděluje
    tyto dvě třídy. Nemusíte nic počítat.
  \item Aniž byste cokoliv počítali, vysvětlete, jak přidání nového bodu $(3,3)$ do třídy A ovlivní
    rozhodovací hranici lineárního SVM.
  \item Vylepšilo by na datech uvedených výše použití RBF jádra klasifikaci? Zakreslete nějaký
    dataset, kde by RBF jádra pomohla.
\end{enumerate}

\medskip
\noindent\textbf{Řešení.} \textit{(V~PDF bez náčrtu řešení; vlastní vzorové řešení, separátor ověřen řešičem SVM (\texttt{sklearn}).)}\par
\begin{enumerate}[nosep]
  \item Lineární SVM hledá oddělující přímku (obecně nadrovinu) s~\emph{maximálním marginem} -- největším odstupem od nejbližších bodů obou tříd. Řeší úlohu kvadratického programování $\min\tfrac12\|\mathbf{w}\|^2$ za podmínek $t_i(\mathbf{w}^\top \mathbf{x}_i+b)\ge 1$. Hranici určují \emph{podpůrné vektory} (body s~nenulovým multiplikátorem); ty leží přesně na okraji marginu, obráceně to platit nemusí.
  \item Nejbližší body třídy A leží na přímce $x_1+x_2=2$ ($(1,1)$ a~$(2,0)$), nejbližší body třídy B na $x_1+x_2=1$ ($(1,0)$ a~$(0,1)$). Maximálně marginová hranice je rovnoběžka uprostřed, $x_1+x_2=1{,}5$: bod $(1,1)$ má od úsečky $(1,0)$--$(0,1)$ (leží v~konvexním obalu třídy B) vzdálenost $1/\sqrt2$, takže žádná oddělující přímka nemá větší margin než polovinu, a~tato přímka ho dosahuje (kanonicky $\mathbf{w}=(2,2)$, $b=-3$, margin $1/(2\sqrt2)\approx0{,}354$). Podpůrné vektory (s~nenulovým multiplikátorem) vrací řešič $(1,1)$ ze třídy A a~$(1,0),(0,1)$ ze třídy B; bod $(2,0)$ leží na marginu také.
  \item Bod $(3,3)$ leží daleko od hranice na správné straně, není tedy podpůrným vektorem -- rozhodovací hranice se \emph{nezmění}.
  \item Data jsou už lineárně separabilní, takže RBF jádro klasifikaci nezlepší (spíš by hrozilo přeučení). RBF pomáhá u~\emph{nelineárně} oddělitelných dat -- např. jedna třída tvoří shluk uvnitř a~druhá prstenec kolem něj (soustředné kružnice).
\end{enumerate}
